\begin{recipedef}
  Name:                & \REC{metis_n_img_std_process} \\
  Purpose:             & Determine the conversion factor between detector counts and physical source flux \\
%  Requirements:        & \REQ{METIS-XXXX} \\
%  Type:                & Calibration | Science | Technical \\
%  Templates:           & \TPL{...} \\
  Primary input:       & \PROD{N_STD_BKG_SUBTRACTED} (N$\times$) \\
  Secondary inputs:    & \EXTCALIB{FLUXSTD_CATALOG} \\
  Matched keywords:    & \FITS{DRS.FILTER} \\
  Parameters:          & \CODE{metis_n_img_std_process.stacking.method}: string\newline Name of the method used to combine the input images (\texttt{add}, \underline{\texttt{average}}, \texttt{median}, \texttt{sigclip}) \\
  Algorithm:           & Call metis\_n\_calculate\_std\_flux to measure flux in input images \\
                       & call hdrl\_resample\_compute to recenter the images \\
                       & call hdrl\_imagelist\_collapse to stack the images \\
                       & call metis\_n\_calculate\_std\_flux on the stacked image to get flux of the star in detector units \\
                       & call metis\_calculate\_std\_fluxcal to calculate the conversion factor to physical units \\
                       & call metis\_calculate\_detection\_limits to compute background noise (std, rms) and compute detection limits \\
  Output data:         & \PROD{FLUXCAL_TAB} \\
                       & \PROD{N_STD_COMBINED} \\
%  Expected accuracies: & \\
  QC1 parameters:      & \QC{QC N STD AIRMASS} \\
                       & \QC{QC N STD APERTURE CNTS} \\
                       & \QC{QC N AREA SENS} \\
                       & \QC{QC N IMG STD BACKGD RMS} \\
                       & \QC{QC N STD ELLIP} \\
                       & \QC{QC N STD FLUXCONV} \\
                       & \QC{QC N STD FWHM} \\
                       & \QC{QC N STD PEAK CNTS} \\
                       & \QC{QC N SENS} \\
                       & \QC{QC N STD STREHL} \\
\end{recipedef}

